\providecommand{\pdfxopt}{a-1b,mathxmp}

%\documentclass[10pt]{article}
%\usepackage{beamerarticle}
\documentclass[ignorenonframetext,slidestop]{beamer}

\input{/home/asreeve/current/tex_style/preamble_102}
\usepackage{cancel}

\setbeamertemplate{enumerate item}[square]
\setbeamertemplate{enumerate subitem}[circle]
\setbeamertemplate{enumerate subsubitem}[square]
\setbeamercolor*{enumerate item}{fg=red,bg=blue}
\setbeamercolor*{enumerate subitem}{fg=blue,bg=red}

\begin{document}

\begin{frame}{Peat as a Resource}
  \begin{enumerate}
  \item Peat is an organic sediment that accumulates in some wetland environments.
  \item Some peat deposits are precursors to coal formation
  \item Peat is used a fuel, as a soil additive, and for cosmetic purposes (e.g.peat baths).
  \end{enumerate}

  In this lab, you will explore the energy value of peat deposits in Maine and work through the calculations to estimate the amount of peat in a Maine peatland.
\end{frame}

\begin{frame}{Calculating Irregular Areas and Volumes}
  \begin{itemize}
  \item How can the areas of irregular objects, like a peatland, be estimate?
  \item complex math and computer methods like digitizing with a computer
  \item simpler hand methods using geometry and trigonometry
    \begin{itemize}
    \item Grid method
      \begin{enumerate}
      \item overlying a grid,
      \item determine grid cell area
      \item counting grid cells in irregular object
      \item multiply cell area by number of cells
      \end{enumerate}
    \item Polygon method
      \begin{enumerate}
      \item divide irregular object into polygons
      \item calculate area of polygons 
      \item sum areas of all the polygons
      \end{enumerate}
    \end{itemize}
  \end{itemize}
\end{frame}

\begin{frame}{Getting the area using a grid}
  
  Using first method: overlay grid on irregular object (graphics program over digital image, graph paper on printed image).

  \begin{minipage}{.63\textwidth}
    \includegraphics[width=\textwidth]{../../figs/peat_area_grid.pdf}
  \end{minipage}
  \begin{minipage}{.35\textwidth}
    \begin{itemize}
    \item Note scale on image. 10 cells span what 100 m on map.
    \item Cell is 100m/10 = 10 m across
    \item one 10 m by 10 m call: $ 10 m \times 10 m = 100 m^2$
    \item  cell area = 100 $m^2$
    \end{itemize}
  \end{minipage}
\end{frame}

\begin{frame}{Getting the area using a grid}
  Area of heart shaped blob in graphic: 1. count cell vertices (where lines cross) inside object, 2. area around the vertices is cell area (100 $m^2$), (as shown).
  
  \begin{minipage}{.63\textwidth}
    \includegraphics[width=\textwidth]{../../figs/peat_area_grid2.pdf}
  \end{minipage}
  \begin{minipage}{.35\textwidth}
    \begin{itemize}
      \item Count number of vertices.
      \item done in small groups along rows, that added together.
      \item 487 total cells in irregular heart-blob.
      \item number of cells multiplied by cell area = total area of object.
        \scriptsize $ 487 \times 100\, m^2 =48700 \, m^2$
    \end{itemize}
  \end{minipage}
\end{frame}

\begin{frame}{Breaking the object into polygons}
  Polygonal method: break up object into simple polygons. Recalling high school mathematics:
  \begin{itemize}
  \item Triangle area $= \frac{1}{2} base \times height$
  \item Rectangle area $= base \times height$
  \end{itemize}
\end{frame}
\begin{frame}{Breaking the object into polygons}
  In image below, irregular shape broken into rectangles and triangle. Using equations for polygons and map scale, can estimate real area object.
  
   \begin{minipage}{.58\textwidth}
    \includegraphics[width=\textwidth]{../../figs/peat_area_poly.pdf}
  \end{minipage}
  \begin{minipage}{.40\textwidth}
    \begin{itemize}
      \item yellow triangle area of 2100 $m^2$
      \item orange triangle area of 400  $m^2$
      \item pink rectangle area of 5600 $m^2$
      \item green rectangle area of 14000  $m^2$
      \end{itemize}
          \end{minipage}
Calculate area for all polygons, sum polygon areas to estimate total  object area.
\end{frame}

\begin{frame}{Estimating peat area and volume}
  \begin{minipage}{.45\textwidth}
    \begin{itemize}
    \item When estimating peatland area, only use 'heath' units
    \item To estimate peat volume, multiply area by average thickness of low ash peat (darkest colored band)
    \item Low ash peat depths shown in 'explanation of cores'. Numbers on map correlate with core number
    \end{itemize}
  \end{minipage}
  \begin{minipage}{.5\textwidth}
    \includegraphics[width=\textwidth]{../../figs/sunkhaze_map1_2.png}
  \end{minipage}
\end{frame}

\begin{frame}{Doing unit conversions}
  This lab includes conversions between units. When converting between measurements that cross at zero \footnotesize (inches and cm for example, 0 inches = 0 cm) \normalsize,
  \begin{enumerate}
  \item multiply by a fraction with equal amount but different units in numerator and denominator. 
  \item arrange fractions so the units cancel, leaving only the desired unit
  \end{enumerate}

  \footnotesize
  \begin{block}{Example: How many cm are in a mile?}
    \begin{enumerate}
    \item identify conversions that will allow stepping from 1 mile to cm.
      \begin{itemize} \footnotesize
      \item 1 mile = 5280 ft.
      \item 1 ft = 12 inches
      \item 1 inch = 2.54 cm
      \end{itemize}
    \item write out equation so that only desired units do not cancel:
    \end{enumerate}
    \begin{equation*}
      \underbrace{1 \; \cancel{mile}}_{starting \; value} \times \frac{5280 \; \bcancel{ft}}{1 \; \cancel{mile}} \times
      \frac{12 \; \xcancel{inches}}{1 \; \bcancel{ft}} \times \frac{2.54 \; cm}{1 \; \xcancel{inch}}
      = 160934.4 \; cm
    \end{equation*}
    or 161000 cm if using significant digits
  \end{block}
\end{frame}

\begin{frame}{Doing unit conversions}
 \footnotesize
 \begin{block}{Example: How much energy is in 27 ft$^3$ of wet peat?}
   \begin{itemize}
   \item 1 wet ft$^3$ peat = 8.7 lbs dried peat (in lab pdf)
   \item 1 lbs dried peat = 9518 BTUs (where did this come from in the lab pdf?)
   \end{itemize}
   \begin{equation*}
     27 \, wet \, ft^3 \times \frac{8.7 \,lbs \,dry}{1 \,wet \,ft^3} \times \frac{9518 \,BTUs}{1 \,lbs \,dry} = 2235778.2 \, BTUs
   \end{equation*}
   or $2.2 \cdot 10^6$  BTUs if using significant figures
 \end{block}
 For those who would like a more in-depth discussion of algebra and unit conversions and a review of significant figures, you may want to read 'CG-1, Significant figures' and 'CG-7, The algebra of unit conversions' available from: \href{https://www.nagt.org/nagt/jge/columns/compgeo.html}{https://www.nagt.org/nagt/jge/columns/compgeo.html}
\end{frame}
\end{document}
